\section{Structural Limitations}
\label{sec:theorems}

\subsection{Nonnegative weights cannot invert a source}

\begin{theorem}[Nonnegative-weight insufficiency]
\label{thm:nonneg}
Suppose sources use inversion-only conventions with full stability
($\rho_s = 1$ for all $s$), and at least one source $j$ uses an inverted
convention ($\sigma_j = -1$). Suppose there exists a message vector
$X_{-j}$ from the remaining sources such that
$\bigl|\sum_{s \neq j} \sigma_s \lambda_s X_s\bigr| < \lambda_j$.
Then no rule in $\mathcal{A}_+$ agrees with the Bayes rule on all message
vectors. Here $\mathcal{A}_+$ is the class of rules
$\hat{y}(X) = \mathrm{sign}(\sum_s w_s X_s)$ with $w_s \geq 0$ fixed per
source at decision time: the weights may be estimated from support data
by any procedure, but they cannot depend on the current message vector
or context.
\end{theorem}

\begin{proof}
Consider the inversion-only regime, where the Bayes log-odds are
$\mathrm{LO}(X) = \sum_s \sigma_s \lambda_s X_s$.
Fix $X_{-j}$ satisfying the hypothesis and form the pair
$X^{(1)} = (X_{-j}, -1)$, $X^{(2)} = (X_{-j}, +1)$, differing only in
source $j$'s message. Writing $S = \sum_{s\neq j}\sigma_s\lambda_s X_s$,
\[
\mathrm{LO}(X^{(1)}) = S + \lambda_j > 0
\quad\text{and}\quad
\mathrm{LO}(X^{(2)}) = S - \lambda_j < 0,
\]
using $\sigma_j = -1$ and $|S| < \lambda_j$. Hence the Bayes decision is
strictly \emph{decreasing} in $X_j$ on this pair: $+1$ then $-1$.
Both log-odds are nonzero (the inequalities are strict), so the
fair-coin tie convention of Section~\ref{sec:formal} never engages on
this pair.

Any rule $U \in \mathcal{A}_+$ computes
$U(X) = \mathrm{sign}(\sum_s w_s X_s)$ with $w_s \geq 0$. Moving from
$X^{(1)}$ to $X^{(2)}$ changes the score by $2w_j \geq 0$, so $U$ is
\emph{non-decreasing} in $X_j$ and cannot reproduce the Bayes flip from
$+1$ to $-1$. Contradiction.
\end{proof}

The argument is a monotonicity obstruction, and it is worth stating
plainly what it rules out: the Bayes rule sometimes has to move
\emph{against} a source's message, and no nonnegative weighting can ever do
that. The hypothesis is sufficient rather than generic. It names concrete
configurations where the failure bites, which is what an impossibility
result is for. The cleanest witness is the single-source case: with
$X_{-j}$ empty, $S = 0 < \lambda_j$ holds trivially, so a lone inverted
source already defeats every nonnegative rule.

\begin{corollary}
In the inversion-only regime with full stability ($\rho_s = 1$), the
signed class $\mathcal{A}$ contains the Bayes rule, attained at
$v_s = \sigma_s \lambda_s$.
\end{corollary}

\subsection{Signed scalars cannot handle context-dependence}

Allowing negative weights fixes inversion. It does not fix the deeper
problem: when a source's convention conditions on context, the source's
Bayes log-odds contribution is itself a function of the context ---
shrinking, vanishing, or flipping sign as stability varies --- and no
single scalar can represent a function.

\begin{theorem}[Scalar insufficiency for context-dependent conventions]
\label{thm:scalar}
Let source $j$ have convention $c$ with flip function $g_c$ (so
$f_c(v,a) = v \oplus g_c(a)$), stability $\rho_j \in (0,1]$ and
reliability $r_j \in (1/2,1)$. Suppose there exist contexts
$a^{(1)} \neq a^{(2)}$ in the support of $P_A$ with
$g_c(a^{(1)}) = 0$ and $g_c(a^{(2)}) = 1$. Write
$\lambda_j = \log\frac{r_j}{1-r_j}$ and let $\varepsilon_a$ be the
induced channel error of Section~\ref{sec:formal}.

(a) [\emph{Log-odds}] Source $j$'s Bayes log-odds contribution for
message $m$ in context $a$ is
$\ell_j(m,a) = (2m-1)\log\frac{1-\varepsilon_a}{\varepsilon_a}$.
In context $a^{(1)}$ the weight is $+\lambda_j$; in context $a^{(2)}$
it is a strictly decreasing function of $\rho_j$, equal to
$+\lambda_j$ at $\rho_j = 0$, to $0$ at $\rho_j = 1/2$, and to
$-\lambda_j$ at $\rho_j = 1$. In particular, whenever $\rho_j > 0$ the
contribution differs across the two contexts, so no
context-independent scalar $w_j$ --- of either sign --- equals source
$j$'s Bayes log-odds contribution in both.

(b) [\emph{Decisions}] If moreover $\rho_j > 1/2$, so that the
$a^{(2)}$-weight $\ell_j^{(2)}$ is negative while the
$a^{(1)}$-weight is positive, fix the other sources' messages and
contexts so that their total Bayes log-odds $S$ satisfies
$|S| < \min(\lambda_j, |\ell_j^{(2)}|)$ (attainable, e.g.\ source $j$
alone, where $S = 0$). Then there exist observations
$x^{(1)} = (X_{-j}, X_j, a^{(1)})$ and
$x^{(2)} = (X_{-j}, X_j, a^{(2)})$ sharing the same centered message
vector $X = (X_{-j}, X_j)$, with $X_j = 2M_j - 1$, on which the Bayes
decisions differ. Consequently no context-independent function of the
message vector alone --- in particular, no scalar-weight rule in
$\mathcal{A}$ or $\mathcal{A}_+$ --- implements the Bayes rule on all
observations. A context-dependent weight $w_j(a)$ is a different
function class, and is not ruled out.
\end{theorem}

\begin{proof}
See Appendix~\ref{app:deferred-proofs}.
\end{proof}

\begin{remark}
Two boundary cases delimit the result. At $\rho_j = 1/2$ a flipped
context contributes exactly zero log-odds: the source is not misread
there, it is silent, and no scalar is contradicted. At
$\rho_j < 1/2$ the flipped-context weight is positive but shrunk, so
part (a) still rules out exact scalar representation of the log-odds,
while decisions need not flip; part (b)'s sign flip needs
$\rho_j > 1/2$, and the full-stability picture --- the weight flipping
from $+\lambda_j$ to $-\lambda_j$ --- is the endpoint $\rho_j = 1$.
Two assumptions deserve to be explicit. First, the receiver
\emph{observes} the context $a$; without that, context-dependent
conventions are trivially undecodable. Second, the convention family
is \emph{known} to the receiver; the theorem is about representational
capacity, not discovery. The obstacle is context-dependence, not sign.
\end{remark}

\subsection{Equal-reliability optimum}

Lemma~\ref{lem:equalrel} quantifies the price of the monotonicity
obstruction: with $p$ honest and $q$ inverted sources at reliability $r$
and full stability, the Bayes accuracy counts all $p+q$ decoded readings,
while the best unsigned rule attains only the honest-only analogue, i.e.\
it optimally discards the inverted sources entirely (statement, proof, and
a numerical illustration are in Appendix~\ref{app:deferred-proofs}).

\subsection{The Bayesian aggregator is consistent}

The results above say what \emph{fixed} rules cannot do. The natural
question is whether the Bayesian aggregator of
Section~\ref{sec:method}, which learns the convention instead of fixing a
weight, actually finds the truth. It does, in the standard sense.

\begin{theorem}[Posterior consistency]
\label{thm:consistency}
Fix a source with true parameters $\theta^* = (c^*, \rho^*, r^*)$ lying on
the finite grid $\mathcal{C} \times \textsc{RhoGrid} \times \textsc{RGrid}$.
Let $[\theta^*]$ be its observational-equivalence class: the grid points
inducing the same distribution over support triples $(M, a, Y)$. Then, as
$K \to \infty$,
\[
P([\theta^*] \mid \mathcal{D}_s) \xrightarrow{a.s.} 1,
\]
and the aggregator's predictive log-odds converge almost surely to the
oracle log-odds (exact Bayes inference with known $\theta^*$).
\end{theorem}

\begin{proof}
See Appendix~\ref{app:deferred-proofs}.
\end{proof}

\begin{remark}
The theorem promises class concentration, not point identification, and
the distinction is load-bearing: on our own grid, distinct parameters can
be observationally identical. Under an inversion-only convention,
$(\rho, r)$ enter the triple distribution only through the agreement
probability $p = \rho(1-r) + (1-\rho)r$, so $(\rho, r) = (1.0, 0.55)$ and
$(0.6, 0.75)$---both on-grid---induce identical support-triple
distributions; under the identity convention, $\rho$ drops out of the
likelihood entirely. What the posterior identifies is the triple
distribution itself, and that is all the predictive needs: the query
likelihood is the same conditional, constant across each equivalence
class, so the predictive log-odds converge to the oracle's regardless.

A caveat on scope: the theorem assumes $\theta^*$ lies exactly on the
grid, while our simulations draw $\rho$ and $r$ continuously. Off-grid,
Berk's theorem~\cite{berk1966} gives concentration on the KL-closest
grid point, but that point is unique only up to predictive equivalence:
where two grid points induce the same support-triple distribution ---
as in the inversion example above --- the posterior may split mass
between them rather than settle on one, and this splitting occurs in
our own runs. The predictive claim is unaffected, because every such
point induces the same query likelihood, so the predictive log-odds
converge to those of the KL-closest predictive distribution. The grid
is fine enough ($\Delta\rho = 0.1$, $\Delta r = 0.05$) that the
projection is close to the truth. We cite Experiment~C
(Section~\ref{sec:experiments}) --- at $K = 50$ the aggregator reaches
$0.8204 \pm 0.0029$ (30 seeds), essentially the expected reliability
ceiling $\mathbb{E}[r] = 0.825$ --- as evidence about the grid's
accuracy only, not as part of the consistency claim.
\end{remark}

\begin{corollary}[Joint aggregator]
\label{cor:joint}
Under the conditional independence of
Section~\ref{sec:formal}, the joint aggregator that sums per-source
predictive log-odds inherits Theorem~\ref{thm:consistency}: if every
source's parameters lie on the grid, the summed predictive log-odds
converge almost surely to the joint oracle log-odds as each
$K_s \to \infty$.
\end{corollary}

\begin{proof}
Each source's predictive log-odds converge a.s.\ to its oracle
log-odds by Theorem~\ref{thm:consistency}; a finite sum of
almost-surely convergent terms converges almost surely, and the sum
of the per-source oracle log-odds is the joint oracle log-odds under
conditional independence.
\end{proof}
